\documentclass[12pt]{article}
\usepackage[a4paper,margin=1in]{geometry}
\usepackage{amsmath, amssymb, booktabs, setspace}
\usepackage{enumitem}
\usepackage{graphicx}
\usepackage{siunitx}

\title{Practice Problem Set 5\\ \vspace{10pt} \large ECON 101 -- Macroeconomics}
\author{}
\vspace{-10pt}
\date{ \vspace{-30pt} Winter 2026}

\begin{document}
\maketitle


%%%%%%%%%%%%%%%%%%%%%%%%%%%%%%%%%%%%%%%%%%%%%%

\begin{enumerate}[label=\textbf{\arabic*.}]

\item Explain the three functions of money:
\begin{itemize}
    \item medium of exchange,
    \item unit of account,
    \item store of value.
\end{itemize}
Provide one real-world example for each.

\item Distinguish between:
\begin{itemize}
    \item M1 and broader measures of money supply,
    \item currency vs. deposits.
\end{itemize}
Why are some assets included in the money supply while others are not?

\item Explain how commercial banks create money. 
What role does the reserve ratio play in this process?

\end{enumerate}



\begin{enumerate}[label=\textbf{\arabic*.}, start=4]

\item \textbf{Deposit Creation}

Suppose:
\begin{itemize}
    \item Initial deposit = \$5,000,
    \item Desired reserve ratio = 0.10.
\end{itemize}

\begin{enumerate}
    \item[(a)] Calculate required reserves.
    \item[(b)] Calculate initial loans.
    \item[(c)] Calculate the simple deposit multiplier.
    \item[(d)] Calculate the maximum increase in deposits.
\end{enumerate}

\item \textbf{Leakages in the Banking System}

Explain how the following reduce the money multiplier:
\begin{itemize}
    \item excess reserves,
    \item currency held by the public.
\end{itemize}

\end{enumerate}


\begin{enumerate}[label=\textbf{\arabic*.}, start=6]

\item Explain how a bank run occurs and how the Bank of Canada and deposit insurance help prevent or limit bank runs.

\end{enumerate}


\vspace{1em}

%%%%%%%%%%%%%%%%%%%%%%%%%%%%%%%%%%%%%%%%%%%%%%



\begin{enumerate}[label=\textbf{\arabic*.}, start=7]

\item What is the Bank of Canada’s inflation target? 
Why is inflation targeting important?

\item Explain the monetary transmission mechanism step-by-step starting from a change in the policy interest rate.

\item Distinguish between:
\begin{itemize}
    \item expansionary monetary policy,
    \item contractionary monetary policy.
\end{itemize}
When would each be used?

\end{enumerate}



\begin{enumerate}[label=\textbf{\arabic*.}, start=10]

\item \textbf{Multiplier and Monetary Policy}

Suppose:
\begin{itemize}
    \item MPC = 0.8,
    \item Investment increases by \$40 billion due to lower interest rates.
\end{itemize}

\begin{enumerate}
    \item[(a)] Calculate the multiplier.
    \item[(b)] Calculate the total change in GDP.
\end{enumerate}

\item \textbf{Interest Rate Effects}

Explain how an increase in interest rates affects:
\begin{itemize}
    \item consumption,
    \item investment,
    \item exchange rate,
    \item net exports.
\end{itemize}

\end{enumerate}


\begin{enumerate}[label=\textbf{\arabic*.}, start=12]

\item Why might monetary policy be less effective during a recession? 
Discuss at least two reasons.

\end{enumerate}


\vspace{1em}

%%%%%%%%%%%%%%%%%%%%%%%%%%%%%%%%%%%%%%%%%%%%%%



\begin{enumerate}[label=\textbf{\arabic*.}, start=13]

\item Distinguish between:
\begin{itemize}
    \item discretionary fiscal policy,
    \item automatic stabilizers.
\end{itemize}
Provide examples of each.

\item Explain how government spending and taxes affect aggregate demand.

\item What is a government budget deficit and surplus? 
How are they related to the business cycle?

\end{enumerate}


\begin{enumerate}[label=\textbf{\arabic*.}, start=16]

\item \textbf{Government Spending Multiplier}

Suppose:
\begin{itemize}
    \item MPC = 0.75,
    \item Government spending increases by \$80 billion.
\end{itemize}

\begin{enumerate}
    \item[(a)] Calculate the multiplier.
    \item[(b)] Calculate the change in GDP.
\end{enumerate}

\item \textbf{Tax Multiplier}

Suppose:
\begin{itemize}
    \item MPC = 0.75,
    \item Taxes decrease by \$60 billion.
\end{itemize}

\begin{enumerate}
    \item[(a)] Calculate the tax multiplier.
    \item[(b)] Calculate the change in GDP.
\end{enumerate}

\end{enumerate}



\begin{enumerate}[label=\textbf{\arabic*.}, start=18]

\item The economy is in a recession:
\begin{itemize}
    \item GDP is below potential,
    \item unemployment is rising.
\end{itemize}

\begin{enumerate}
    \item[(a)] Use the AE model to explain the recession.
    \item[(b)] Use the AD--AS model to illustrate the situation.
    \item[(c)] Propose one fiscal policy and one monetary policy. 
    Explain how each affects GDP.
\end{enumerate}

\end{enumerate}


\begin{enumerate}[label=\textbf{\arabic*.}, start=19]

\item Suppose government increases spending by \$100 billion but finances it through borrowing.

Explain:
\begin{itemize}
    \item how interest rates change,
    \item how private investment is affected,
    \item why the actual increase in GDP may be smaller than predicted.
\end{itemize}

\end{enumerate}

\hrule
\vspace{1em}

\begin{center}
\textit{End of Problem Set}
\end{center}

\end{document}